% ============================================================================
%  CHAPITRE 8 -- SIMULATEUR INTERACTIF
% ============================================================================
\chapter{Simulateur interactif}
\label{chap:simulateur}

Les courbes du chapitre~\ref{chap:resultats} donnent des moyennes sur des milliers d'images, mais ne montrent
pas ce qui se passe \emph{à l'intérieur} d'une transmission. Le simulateur la déroule étape par étape.

\begin{lstlisting}[language=bash, caption={Lancement}, label={lst:explore}]
python scripts/explore.py          # ouvre http://127.0.0.1:8077/
\end{lstlisting}

\section{Principe}
\label{sec:sim_principe}

Aucune étape n'est rejouée ni pré-enregistrée : chaque étape appelle le code du projet
(\srcfile{semcom.channel}, \srcfile{semcom.digital}, \srcfile{semcom.models}, \srcfile{semcom.metrics}). Ce
que la page affiche est donc ce que \srcfile{scripts/evaluate.py} mesure. Le simulateur observe la chaîne, il
ne la réimplémente pas.

\begin{table}[H]
\centering
\caption{Modules du simulateur}
\label{tab:sim_modules}
\small
\begin{tabularx}{\textwidth}{@{}p{3.3cm}>{\raggedright\arraybackslash}X@{}}
\toprule
\textbf{Module} & \textbf{Rôle} \\
\midrule
\srcfile{explorer/trace.py} & exécute une transmission et enregistre chaque étape \\
\srcfile{explorer/catalog.py} & systèmes, images et options ; chargement des modèles \\
\srcfile{explorer/encode.py} & conversion des tenseurs en objets affichables \\
\srcfile{explorer/server.py} & serveur HTTP et API JSON \\
\srcfile{explorer/static/} & carte animée, lecteur d'étapes, inspecteur, courbes \\
\bottomrule
\end{tabularx}
\end{table}

Le serveur repose sur \srcfile{http.server} et l'interface sur du HTML, CSS et JavaScript sans bibliothèque :
aucune dépendance n'est ajoutée au projet.

\section{Ce que l'on observe}
\label{sec:sim_vues}

Émetteur, canal et récepteur forment une carte sur laquelle progresse l'étape courante. Chaque étape affiche
la grandeur réellement présente à cet endroit :

\begin{itemize}
  \item le \textbf{tenseur latent} en pile 3D, canaux transmis en couleur et canaux tronqués en gris ;
  \item les \textbf{symboles I/Q} avant et après bruit : nuage continu pour le JSCC, constellation discrète
        pour la QAM ;
  \item les \textbf{bits codés}, bits faux marqués, avec le taux d'erreur avant et après Viterbi ;
  \item les \textbf{LLR} signés ; l'\textbf{image reçue} ; le \textbf{verdict des juges}.
\end{itemize}

Quatre onglets : \emph{Explorer}, \emph{Comparer} (plusieurs systèmes sur la même image, calculés en direct),
\emph{Résultats mesurés} (courbes lues dans les CSV) et \emph{Concepts}. L'état tient dans l'URL, donc une
étape précise se partage par lien.

\section{Apports}
\label{sec:sim_apport}

\paragraph{Pédagogique.} L'onglet \emph{Comparer} à \SI{2}{\decibel} montre la thèse du projet en une vue : la
chaîne séparée rend une image grise, le JSCC syntaxique une image reconnaissable mais mal classée, le JSCC
sémantique une image correctement classée.

\paragraph{Vérification.} Les grandeurs intermédiaires rendent observables des comportements que les moyennes
masquent : le budget binaire insuffisant pour un seul canal latent aux MCS robustes
(section~\ref{sec:codage_source}), ou le fait que très en dessous du seuil, Viterbi peut produire un taux
d'erreur \emph{supérieur} au taux brut en s'engageant sur un mauvais chemin du treillis. Deux conséquences
correctes de la théorie, invisibles sur une courbe de PSNR.
